\documentclass[11pt]{elsarticle}

\usepackage{graphicx}
\graphicspath{{./}{../figures/}{figures/}}
\usepackage{booktabs}
\usepackage{amsmath,amssymb}
\usepackage{longtable}
\usepackage{url}
\usepackage{float}
\usepackage{placeins}
\usepackage[margin=2.4cm]{geometry}

\renewcommand{\thetable}{S\arabic{table}}
\renewcommand{\thefigure}{S\arabic{figure}}
\renewcommand{\thesection}{S\arabic{section}}

\begin{document}

\begin{center}
{\Large \textbf{Supplementary Material}}\\[6pt]
{\large Directional transfer of chemical-gene effects from rodents to humans:
an independence-guarded meta-analysis of the Comparative Toxicogenomics
Database}
\end{center}

\vspace{4pt}

\noindent This file contains the supplementary tables referenced in the main
text. Unless stated otherwise, the analytic population is the primary set:
chemical-gene-endpoint effect keys measured in both a human and a rodent
context, with an unambiguous consensus direction in each species, and with
disjoint supporting PubMed identifiers across species (n=40{,}779). Agreement is
the probability that the human consensus direction equals the rodent consensus
direction. Point estimates carry Wilson 95\% intervals; interval estimates
labelled ``bootstrap'' resample chemicals (1{,}000 replications) to respect the
nesting of effect keys within chemicals. All six figures referenced in this
supplement appear in the main text (Figures 1-6).

\vspace{6pt}

\noindent \textbf{Contents.} The four headline comparisons (overall independent
agreement, all pairs, shared-publication pairs, and the mouse-versus-rat
within-rodent benchmark), with both Wilson and cluster-bootstrap intervals, are
given in Table 1 of the main text and are not repeated here. The tables below
provide the stratified and mechanistic detail that supports the main text.

\begin{itemize}\itemsep2pt
\item Table S1. Agreement by molecular endpoint.
\item Table S2. Agreement by replication depth.
\item Table S3. Agreement by chemical class.
\item Table S4. Within-species directional inconsistency.
\item Table S5. Variance decomposition of the agreement outcome (intraclass correlations).
\item Table S6. Predictive modelling: grouped cross-validation, permutation importance, and leave-one-class-out generalisation.
\item Table S7. External conservation anchor: agreement by protein-identity quartile.
\item Table S8. External conservation anchor: cluster-robust trend, one-to-one orthology, and incremental predictive value.
\item Table S9. Catalogue of independent cross-species directional conflicts.
\end{itemize}

\clearpage
%=========================================================================
\section{Stratified reproducibility}
%=========================================================================

\begin{table}[H]
\small
\centering
\caption{Independent human-rodent directional agreement by molecular endpoint
(primary set), ordered by sample size. Endpoints with fewer than 30 pairs are
omitted. Methylation agreement is indistinguishable from its chance baseline.}
\label{tab:s1}
\begin{tabular}{lcccc}
\toprule
Molecular endpoint & n & Agreement & Bootstrap 95\% CI & Chance \\
\midrule
Expression & 36{,}986 & 60.8\% & 58.3-63.6 & 50.4\% \\
Activity & 1{,}106 & 89.3\% & 87.2-91.2 & 50.8\% \\
Methylation & 1{,}102 & 41.9\% & 37.1-64.4 & 40.9\% \\
Phosphorylation & 744 & 81.5\% & 77.5-85.0 & 64.0\% \\
Secretion & 286 & 85.7\% & 81.1-90.6 & 75.1\% \\
Cleavage & 182 & 96.7\% & 94.2-98.9 & 95.7\% \\
Response to substance & 144 & 75.0\% & 67.3-81.9 & 50.3\% \\
Metabolic processing & 55 & 98.2\% & 94.2-100 & 98.2\% \\
Abundance & 53 & 84.9\% & 73.2-95.4 & 71.2\% \\
Uptake & 44 & 95.5\% & 89.1-100 & 95.5\% \\
Transport & 40 & 92.5\% & 80.5-100 & 92.7\% \\
\bottomrule
\end{tabular}
\end{table}

\begin{table}[H]
\small
\centering
\caption{Independent directional agreement by replication depth, the minimum
number of supporting publications per species. Agreement rises monotonically to
determinism.}
\label{tab:s2}
\begin{tabular}{lccc}
\toprule
Minimum supporting reports per species & n & Agreement & Bootstrap 95\% CI \\
\midrule
1 & 38{,}135 & 60.7\% & 57.6-63.7 \\
2 & 1{,}980 & 76.5\% & 70.1-84.9 \\
3 & 374 & 90.6\% & 85.7-96.8 \\
4 & 135 & 95.6\% & 92.6-98.4 \\
5 & 60 & 100\% & 100-100 \\
$\geq$6 & 95 & 100\% & 100-100 \\
\bottomrule
\end{tabular}
\end{table}

\begin{table}[H]
\small
\centering
\caption{Independent directional agreement by coarse chemical class (MeSH
chemical tree). The spread across classes is small, consistent with the
near-zero chemical-class variance component in Table S5.}
\label{tab:s3}
\begin{tabular}{lccc}
\toprule
Chemical class & n & Agreement & Bootstrap 95\% CI \\
\midrule
Heterocyclic organic & 4{,}637 & 68.6\% & 63.3-74.3 \\
Other & 15{,}263 & 63.3\% & 60.3-66.4 \\
Solvents / industrial & 2{,}673 & 61.7\% & 58.4-69.6 \\
Organic pollutants & 13{,}853 & 59.3\% & 54.2-67.8 \\
Metals / inorganic & 4{,}336 & 59.1\% & 53.5-70.1 \\
\bottomrule
\end{tabular}
\end{table}

\clearpage
%=========================================================================
\section{Decomposition of non-agreement}
%=========================================================================

\begin{table}[H]
\small
\centering
\caption{Within-species directional inconsistency. For keys curated more than
once within a single species, the fraction that were internally ambiguous (the
same chemical-gene-endpoint curated both up and down). This is a within-species
reproducibility deficit that is present before any cross-species comparison.}
\label{tab:s4}
\begin{tabular}{lccc}
\toprule
Species & Keys curated & Keys curated $>$1 time & Internally ambiguous \\
\midrule
Human & 580{,}801 & 37{,}780 & 58.3\% \\
Mouse & 380{,}904 & 30{,}685 & 54.5\% \\
Rat & 299{,}647 & 22{,}840 & 65.9\% \\
\bottomrule
\end{tabular}
\end{table}

\begin{table}[H]
\small
\centering
\caption{Variance decomposition of the binary agreement outcome (primary set),
as one-way intraclass correlation coefficients (ICC) by grouping factor. The ICC
is the share of agreement variance attributable to differences between groups.
Transferability localises to the specific chemical-gene pair; chemical class
explains almost none of the variance.}
\label{tab:s5}
\begin{tabular}{lcc}
\toprule
Grouping factor & Groups & ICC \\
\midrule
Chemical-gene pair & 40{,}200 & 0.442 \\
Molecular endpoint & 17 & 0.117 \\
Gene & 11{,}949 & 0.068 \\
Chemical & 1{,}477 & 0.046 \\
Chemical class & 7 & 0.005 \\
\bottomrule
\end{tabular}
\end{table}

\noindent Human-specific divergence (main text, Figure 4A): human-versus-rodent
directional disagreement was 38.0\% against a mouse-versus-rat disagreement
ceiling of 33.8\%, an excess of 4.2 percentage points (cluster-bootstrap 95\% CI
$-$1.1 to 10.9), not distinguishable from zero.

\clearpage
%=========================================================================
\section{Predictive modelling}
%=========================================================================

\begin{table}[H]
\small
\centering
\caption{Predicting directional transfer for unseen chemicals under grouped
(by chemical) five-fold cross-validation, so that no chemical appears in both
training and test folds. AUC, area under the ROC curve; AP, average precision;
Brier, mean squared error of the predicted probability. The prevalence baseline
predicts the overall agreement rate (62.0\%) for every effect.}
\label{tab:s6}
\begin{tabular}{lcccc}
\toprule
Model & AUC (mean $\pm$ SD) & AP & Brier \\
\midrule
Gradient boosting & 0.583 $\pm$ 0.025 & 0.725 & 0.231 \\
Logistic regression & 0.562 $\pm$ 0.014 & 0.712 & 0.230 \\
Prevalence baseline & 0.500 $\pm$ 0.000 & 0.620 & 0.236 \\
\bottomrule
\end{tabular}

\vspace{8pt}

\begin{tabular}{lc}
\multicolumn{2}{l}{\textit{Permutation importance (gradient boosting; AUC drop when feature permuted)}} \\
\toprule
Feature & AUC drop \\
\midrule
Rodent consensus direction & 0.056 \\
Gene promiscuity (log) & 0.037 \\
Molecular endpoint & 0.021 \\
Chemical promiscuity (log) & 0.016 \\
Rodent supporting publications & 0.013 \\
Rodent support (vote count) & 0.000 \\
Gene pathway degree (log) & $-$0.001 \\
Chemical class & $-$0.009 \\
\bottomrule
\end{tabular}

\vspace{8pt}

\begin{tabular}{lccc}
\multicolumn{4}{l}{\textit{Leave-one-chemical-class-out generalisation (gradient boosting)}} \\
\toprule
Held-out class & n & AUC & Observed agreement \\
\midrule
Solvents / industrial & 2{,}673 & 0.610 & 61.7\% \\
Heterocyclic organic & 4{,}637 & 0.604 & 68.6\% \\
Metals / inorganic & 4{,}336 & 0.571 & 59.1\% \\
Other & 15{,}263 & 0.564 & 63.3\% \\
Organic pollutants & 13{,}853 & 0.511 & 59.3\% \\
\bottomrule
\end{tabular}
\end{table}

\clearpage
%=========================================================================
\section{External conservation anchor}
%=========================================================================

\noindent The anchor was pre-registered before modelling: human-rodent protein
percent identity (Ensembl Compara, per-gene mean across mouse and rat) and
one-to-one orthology (Ensembl and NCBI \texttt{gene\_orthologs}), all external to
CTD. Coverage of the primary set was near-complete: 39{,}933 of 40{,}779 keys
(97.9\%) carried an Ensembl identity value (median 87.5\%).

\begin{table}[H]
\small
\centering
\caption{Independent directional agreement by quartile of human-rodent protein
identity. Agreement is flat across the full identity range.}
\label{tab:s7}
\begin{tabular}{lcccc}
\toprule
Identity quartile & Median identity & n & Agreement & Bootstrap 95\% CI \\
\midrule
Q1 (lowest) & 70.4\% & 9{,}987 & 63.0\% & 60.6-66.2 \\
Q2 & 83.4\% & 9{,}987 & 61.8\% & 58.5-65.2 \\
Q3 & 90.5\% & 9{,}976 & 62.9\% & 59.9-66.1 \\
Q4 (highest) & 96.7\% & 9{,}983 & 60.5\% & 56.5-64.4 \\
\bottomrule
\end{tabular}
\end{table}

\begin{table}[H]
\small
\centering
\caption{External-anchor association with directional transfer. The trend is a
chemical-cluster-robust logistic regression of agreement on protein identity;
the orthology rows are odds ratios contrasting one-to-one orthologs against all
others. The incremental rows compare the Aim 3 gradient-boosting model with and
without the two conservation features under the same grouped cross-validation,
restricted to keys with a conservation value. All anchors are null.}
\label{tab:s8}
\begin{tabular}{lcc}
\toprule
Anchor / model & Estimate (95\% CI) & $p$ \\
\midrule
Protein identity, per +10 percentage points & OR 0.97 (0.94-1.00) & 0.07 \\
Ensembl one-to-one orthology & OR 1.02 (0.96-1.09) & 0.49 \\
NCBI one-to-one orthology & OR 0.92 (0.78-1.09) & 0.33 \\
\midrule
CTD-internal model & AUC 0.574 (SD 0.030) & -- \\
CTD-internal + conservation & AUC 0.576 (SD 0.023) & -- \\
Difference (paired across folds) & $\Delta$AUC +0.002 (SD 0.009) & -- \\
\bottomrule
\end{tabular}
\end{table}

\clearpage
%=========================================================================
\section{Catalogue of independent cross-species conflicts}
%=========================================================================

\begin{table}[H]
\small
\centering
\caption{Effect keys that are independently and repeatedly curated in
\emph{opposite} directions in humans and rodents (disjoint supporting
publications, unambiguous consensus in each species). These are candidate true
species divergences or curation errors for targeted review. The direction column
gives the human then rodent consensus; ``n PMIDs'' gives supporting publications
on each side. A representative selection is shown; the complete catalogue is
distributed with the analysis code.}
\label{tab:s9}
\begin{tabular}{llllcc}
\toprule
Chemical & Gene & Endpoint & Human / rodent & Human PMIDs & Rodent PMIDs \\
\midrule
Bisphenol A & RPL5 & expression & decrease / increase & 4 & 5 \\
Paraquat & GCLC & expression & decrease / increase & 2 & 7 \\
Dioxin (TCDD) & ALPL & expression & increase / decrease & 2 & 7 \\
Dioxin (TCDD) & IL33 & expression & increase / decrease & 2 & 7 \\
Bisphenol A & ARL3 & expression & increase / decrease & 2 & 6 \\
Bisphenol A & RPL22L1 & expression & decrease / increase & 3 & 5 \\
Bisphenol A & RPL9 & expression & decrease / increase & 2 & 6 \\
Resveratrol & IGF1R & expression & decrease / increase & 4 & 4 \\
Silicon dioxide & LCN2 & expression & decrease / increase & 3 & 6 \\
Bisphenol A & RPL13A & expression & decrease / increase & 3 & 5 \\
Bisphenol A & RPL7 & expression & decrease / increase & 2 & 4 \\
Cadmium chloride & CTNNB1 & expression & increase / decrease & 2 & 5 \\
Dioxin (TCDD) & IGFBP6 & expression & decrease / increase & 3 & 4 \\
MPP$^{+}$ & LDHA & expression & increase / decrease & 1 & 5 \\
4-Nonylphenol & BAX & expression & decrease / increase & 2 & 4 \\
Arsenite & IL6 & expression & decrease / increase & 3 & 4 \\
Bisphenol A & MYO18A & expression & increase / decrease & 2 & 4 \\
Bisphenol A & TBRG4 & expression & decrease / increase & 3 & 3 \\
\bottomrule
\end{tabular}
\end{table}

\noindent The catalogue is enriched for intensively studied environmental
chemicals (bisphenol A, dioxin, paraquat, cadmium) acting on ribosomal and
stress-response genes, the regime in which high-throughput expression profiling
produces the most context-dependent directional calls.

\end{document}
